研究工作按主要研究内容推进：依赖关系建模、PDL 跨边界重叠、全图边处理搜索，以及框架实现与端到端验证。阶段安排与可验收成果如下。

\begin{table}[htbp]
  \centering
  \small
  \caption{研究进度安排与阶段成果}
  \label{tab:schedule}
  \begin{tabularx}{\textwidth}{p{0.18\textwidth}p{0.47\textwidth}X}
    \toprule
    时间 & 主要工作 & 阶段成果 \\
    \midrule
    2026.07--2026.08
      & 完成算子图到 task graph 的 lowering；实现相邻算子边的依赖抽象与合法候选空间；在真实计算图上校验依赖还原
      & 依赖分析 IR、候选生成接口与依赖报告 \\
    2026.09--2026.10
      & 实现依赖感知的 PDL 跨边界变换；识别重叠工作并确定发起时机与执行形态；相对 grid completion 完成边级重叠验证
      & PDL 变换原型、合法性检查与边级重叠结果 \\
    2026.11--2026.12
      & 实现边处理与 grid 分组的联合搜索及成本估计；在 decoder layer 上比较固定方案、逐边贪心与全图搜索
      & 全图计划搜索器及 decoder-layer 时延结果 \\
    2027.01--2027.02
      & 贯通计划生成、CUDA 代码生成与运行时执行；接入完整 decode step 与连续 token 生成，完成正确性通路
      & 可运行框架原型与端到端正确性结果 \\
    2027.03--2027.04
      & 在真实推理 workload 上完成精度与性能测试；扫描 batch、context 与模型 shape，并做依赖/PDL/搜索消融
      & 跨图时延、TPOT/P90/P99、搜索开销与消融结果 \\
    2027.05--2027.06
      & 完善系统原型与复现材料，完成论文撰写
      & 依赖感知融合框架、可复现实验包、投稿稿件和学位论文 \\
    \bottomrule
  \end{tabularx}
\end{table}
